% 第 5 章 系统实现与功能展示 (目标 3 页)
% 借鉴国二报告"功能展示"章: 让评委先看到系统真实存在
% 截图前置: 需先完成高危缺口修复(G2/G3/G5/G1)与部署一键化, 见章节规划-v2.md 第六节
\section{系统实现与功能展示}

\subsection{部署架构与一键复现}

系统以三个进程构成完整演示链路:工具模拟服务、防护代理、评测运行器。仓库提供一键启动脚本,完成依赖检查、服务拉起、样本跑通、结果汇总四步,单次全量评测约 30 分钟。部署指南见附录 C。

% TODO: 图 5 截图: 终端一键启动脚本运行成功画面(三个服务就绪+样本进度条)
\placeholderfig{一键启动脚本运行成功(三服务就绪)}{一键部署与启动}

\subsection{真实智能体端到端演示}

在真实 OpenClaw 类智能体上加载 50 类政务技能,智能体可正常完成公文起草、数据查询、会议安排等日常办公任务。图~\ref{fig:demo-normal} 展示智能体正常办公的交互界面。

% TODO: 图 6 截图: OpenClaw 对话界面, 用户发起正常办公请求, 智能体正常完成
\placeholderfig{OpenClaw 正常办公对话(加载 50 类政务 Skill)}{防护系统下的智能体正常办公}
\label{fig:demo-normal}

\subsection{拦截效果展示}

以数据外泄类攻击为例。用户以业务话术诱导智能体将敏感数据外发,智能体生成外发邮件的工具调用。防护代理在响应解析阶段捕获该调用,规则层命中外发地址关键字,调用被阻断,智能体改而输出合规提示。图~\ref{fig:demo-block} 展示拦截瞬间的智能体回复。

% TODO: 图 7 截图: 攻击 case 触发, 智能体回复被代理改写为阻断提示
\placeholderfig{攻击触发与拦截(阻断提示替换外发调用)}{攻击触发与实时拦截}
\label{fig:demo-block}

\subsection{审计日志展示}

每次拦截在审计日志中留下完整记录,含命中规则编号、评审分数、原始调用参数、最终动作。图~\ref{fig:demo-audit} 展示终端检索审计日志的画面。

% TODO: 图 8 截图: 终端 cat audit_log.jsonl 展示结构化审计记录(含规则ID/评分/trace)
\placeholderfig{审计日志检索(规则编号+评分+调用轨迹)}{审计溯源日志}
\label{fig:demo-audit}
